% Clam AntiVirus: User Manual
% (c) 2002, 2003 Tomasz Kojm <zolw@konarski.edu.pl>
% Version 0.2x corrected by Dennis Leeuw <dleeuw@made-it.com>
\documentclass[a4paper,titlepage,12pt]{article}
\usepackage{amssymb}
\usepackage{pslatex}
\usepackage[dvips]{graphicx}
\usepackage{wrapfig}
\usepackage{url}
%\usepackage{fancyhdr}
\usepackage{fancybox}
\usepackage{titlesec}
\cornersize{.4}
%\fancyhf{}
\headheight 14pt

\usepackage{color}
\definecolor{grey1}{gray}{0.8}
\definecolor{grey2}{gray}{0.3}

% Based on Antonina Liedtke's article in Linux+ 6/2003
\def\greyp{%
    \unitlength=1mm%
    \begin{picture}(0,0)
	\put(0,-1.5){\textcolor{grey1}{\rule{12.9cm}{5.3mm}}\textcolor{grey2}%
	{\rule{9mm}{5.3mm}}\hss}
    \end{picture}
}

\usepackage{fancyhdr}
\pagestyle{fancy}
\fancyhead{}
\fancyfoot{}
\renewcommand{\headrulewidth}{0pt}
\fancyhead[RO]{\textbf{\sffamily{{\textcolor{white}{\thepage}}~}}}
\fancyhead[RE]{\footnotesize{\nouppercase{\rightmark~}}}
\fancyhead[LO]{\footnotesize{\greyp{\nouppercase{\leftmark}}}}


%%%%%%%%%%%%%%%%%%%%%%%%%%%%


%\fancyhead[L,RO]{\small\bfseries\thepage}
%\fancyhead[LO]{\small\bfseries{CLAM ANTIVIRUS: USER MANUAL}}
\date{}

\newcommand{\pl}{\vspace{.3cm}}
\newcommand{\rc}[2]{\textbf{#1: } #2\\[4pt]}
\newcommand{\up}[2]{\textbf{--#1: } #2\\[4pt]}
\newcommand{\email}[1]{\texttt{#1}}
\newcommand{\cons}[1]{\vspace{2mm} \noindent \ovalbox {\sffamily #1}
		      \vspace{2mm}}

\begin{document}
    \setcounter{page}{0}

    \pagestyle{empty}
    %\includegraphics[width=350pt]{clam.eps}
    \begin{center}
	\huge Clam AntiVirus: User Manual \\
	\vspace{1cm}
	\LARGE version 0.60 \\
	\vspace{5cm}
	\LARGE Tomasz Kojm
    \end{center}

    \newpage
    \pagestyle{fancy}
    \tableofcontents
    \newpage


    \section{Introduction}


    Clam AntiVirus is an anti-virus toolkit for UNIX. The main purpose of this
    software is the integration with mail servers (attachment scanning).
    The package provides a flexible and scalable multi-threaded daemon, a
    command line scanner, and a tool for automatic updating via Internet. The
    programs are based on a shared library distributed with the Clam AntiVirus
    package, which you can use with your own software. The virus database
    is based on the virus database from OpenAntiVirus.org, but contains
    additional signatures (including signatures for popular polymorphic viruses,
    too) and is \textbf{kept up to date}.

    \subsection{Features}

    \begin{itemize}
	\item{GNU GPL v2 license}
	\item{POSIX compliant, portable}
	\item{Secure}
	\item{Very fast}
	\item{Multi-threaded}
	\item{User friendly}
	\item{On-access scanning (Linux only)}
	\item{Detects over 7000 viruses, worms and trojans}
	\item{Supports compressed files and archives}
	\item{Built-in support for RAR (2.0), Zip, Gzip, Bzip2}
    \end{itemize}

    \subsection{Mailing lists}
    There are four mailing lists available:
    \begin{itemize}
    \item \textbf{announce@clamav.elektrapro.com} - info about new versions\\
    (including debian package releases), moderated\footnote{That means, the
    subscribers are not allowed to write into the mailing list}.
    \item \textbf{users@clamav.elektrapro.com} - user questions
    \item \textbf{devel@clamav.elektrapro.com} - developement
    \item \textbf{virusdb@clamav.elektrapro.com} - database update information
    \end{itemize}

    \noindent You can subscribe by sending an empty email to\\
    listname-subscribe@clamav.elektrapro.com, or via www at\\
    \indent \url{http://clamav.elektrapro.com/ml}\\
    After subscribing you must reply to a special message sent at your address.

    \noindent Mailing lists are archived at: \\
    \indent \url{http://archive.elektrapro.com/clamav.elektrapro.com/users/}\\
    \indent \url{http://archive.elektrapro.com/clamav.elektrapro.com/devel/}\\


    \subsection{Virus submitting}
    If you have a virus that is not detected by ClamAV with the latest database,
    please send it (as a normal attachment) to:
    \begin{center}
	\email{virus@clamav.elektrapro.com}
    \end{center}
    If your system doesn't allow you to send infected files, please
    archive the virus sample into a zip archive with password: \emph{virus}



    \section{Installation}

    \subsection{Requirements}
    You will need the \emph{zlib} and \emph{zlib-devel} packages and the
    \emph{gcc} compiler (both 2.9x and 3.x are supported). You can install
    the \emph{bzip2} library (and its development files) to get bzip2
    support, but this is not required.

    \subsection{Supported platforms}
    Clam AntiVirus is prepared for the installation on the following
    operating systems / architectures (tested platforms in brackets):
    \begin{itemize}
	\item{GNU/Linux 2.2/2.4 (All flavours, Intel/SPARC/Alpha/zSeries/S/390)}
	\item{Solaris 2.6/7/8/9 (Intel/SPARC)}
	\item{FreeBSD 4.5/6/7 5.0 (Intel/Alpha)}
	\item{OpenBSD 3.0/1/2 (Intel)}
	\item{AIX 4.1/4.2/4.3/5.1 (RISC 6000)}
	\item{HPUX 11.0}
	\item{SCO UNIX}
	\item{Mac OS X}
	\item{BeOS}
	\item{Cobalt MIPS boxes (RAQ1, RAQ2, QUBE2)}
	\item{Windows/Cygwin}
    \end{itemize}

    Some features may not be available with your operating system. If you have
    run Clam AntiVirus on the system not listed above, please let us know.

    \subsection{Actual versions}
    Clam AntiVirus can be obtained from: \\[4pt]
    \begin{center}
	\url{http://clamav.elektrapro.com}\pl
    \end{center}

    \noindent The site is sponsored by ElektraPro.com

    \subsection{Binary packages}
    There are high quality \emph{deb} and \emph{rpm} packages available
    for Linux. The Debian package is maintained by Magnus Ekdahl and you will
    find it on debian mirrors, \url{http://www.debian.org}.
    The RPM package is maintained by Arkadiusz Miskiewicz and is distributed
    with Polish(ed) Linux Distribution (\url{ftp://ftp.pld.org.pl}). There is
    also the RPM package for Mandrake available, it's maintained by Oden
    Eriksson and can be found on Mandrake mirrors.
    The binary packages for AIX are available in AIX PDSLIB, UCLA
    \url{http://aixpdslib.seas.ucla.edu/packages/clamav.html}. The official
    FreeBSD port is maintained by Masahiro Teramoto. The unofficial port
    for OpenBSD (maintained by Flinn Mueller) is available at:\\
    \url{http://www.activeintra.net/openbsd/article.php?id=5}.



    \subsection{Installation}

    \noindent
    Please read the README file in the current version, because it probably
    contains some important release notes.
    If you are installing Clam AV for the first time, you have to add a new user
    and group to your system - \emph{clamav}: \footnote{Cygwin note:
    If you don't have /etc/passwd, you don't need the \emph{clamav} user/group.}

    \begin{verbatim}
	# groupadd clamav
	# useradd -g clamav -s /bin/false -c "Clam AntiVirus" clamav
    \end{verbatim}

    The above method works on Linux and Solaris, if you don't have
    \emph{groupadd, useradd} please consult your system manual - the section
    about creating new users and groups.
    If you are not a system administrator or won't be using \textbf{clamscan}
    in superuser mode, you may omit this step with the option
    \emph{--disable-clamav} passed to the \emph{configure} script:
    \begin{verbatim}
	$ ./configure --disable-clamav 
    \end{verbatim}
    This disables test for the \emph{clamav} user and group. \textbf{clamscan}
    still requires \emph{clamav} for superuser mode. Please don't set a password
    on this account, just assure it's locked with "\textbf{!}" in
    \emph{/etc/passwd} or \emph{/etc/shadow}. It must be a normal, unprivileged
    user. Don't add it to any supplementary groups.\pl

    \noindent
    After you have created the clamav user/group, extract the archive:
    \begin{verbatim}
	$ zcat clamav-x.yz.tar.gz | tar xvf -
	$ cd clamav-x.yz
    \end{verbatim}
    Assuming you want the configuration file installed in /etc, configure
    the package as follows:
    \begin{verbatim}
	$ ./configure --sysconfdir=/etc
    \end{verbatim}
    Currently \emph{gcc} is required for the compilation. Support for other
    compilers will be added in a near future.
    \begin{verbatim}
	$ make
	$ su -c "make install"
    \end{verbatim}
    In the last step the software is installed in the /usr/local directory
    and the config file in /etc.
    \textbf{WARNING: Never set SUID/SGID bit on Clam AntiVirus programs.}

    \subsection{Configuration}
    If you are going to use the daemon, you need to configure it.

    \begin{verbatim}
	$ clamd
	ERROR: Please edit the example config file
	       /etc/clamav.conf.
    \end{verbatim}

    Now you know, where the configuration file is located ;). The format and
    options of this file are fully described in the \emph{clamav.conf(5)}
    manual. clamd configuration is rather easy, the config file is well
    commented. Remember, you must remove the "Example" directive.

    Another feature of clamd is on-access scanning based on the Dazuko module,
    available from \url{http://dazuko.org}. \textbf{This is not required to
    run clamd, furthermore you shouldn't run Dazuko on production systems}.
    A special thread in clamd responsible for the communication with Dazuko
    is called "Clamuko" (it's due to the funny name of Dazuko - I don't know
    what Clamuko means). Clamuko is supported on Linux 2.2 and 2.4 only.
    Dazuko instalation:


    \begin{verbatim}
	$ tar zxpvf dazuko-a.b.c.tar.gz
	$ cd dazuko-a.b.c
	$ make dazuko
	or
	$ make dazuko-smp (for smp kernels)
	$ su
	# insmod dazuko.o
	# cp dazuko.o /lib/modules/`uname -r`/misc
	# depmod -a
    \end{verbatim}

    Depending on your Linux distribution you have to add "dazuko" entry to\\
    \emph{/etc/modules} or the following line:
    \begin{verbatim}
	modprobe dazuko
    \end{verbatim}
    to some startup file to load dazuko at the boot time. You must also create
    the \emph{/dev/dazuko} device:
    \begin{verbatim}
	$ cat /proc/devices | grep dazuko
	254 dazuko
	$ su -c "mknod -m 600 /dev/dazuko c 254 0"
    \end{verbatim}
    Now you must configure Clamuko in \emph{clamav.conf}. Please check
    \ref{clamuko} section.


    \subsection{Testing}

    OK. Let's do some tests. Try to scan the source directory recursively:
    \begin{verbatim}
	$ clamscan -r -l scan.txt clamav-x.yz
    \end{verbatim}
    It should find the viruses in the clamav-x.yz/test directory. You may
    check it in the created log - scan.txt. \textbf{You will find more about
    clamscan options in the clamscan(1) manual. \footnote{Please run \emph{man
    clamscan}}}
    To test clamd first start it and then use clamdscan (you can also connect
    directly to clamd and run the SCAN command):
    \begin{verbatim}
	$ clamdscan -l scan.txt clamav-x.yz
    \end{verbatim}

    \subsection{FreshClam: Setting up auto-updating}
    The \emph{freshclam} utility is the default database updater for Clam
    AntiVirus. It works in two modes:
    \begin{itemize}
	\item interactive - from command line
	\item as a daemon - works alone, silently
    \end{itemize}
    When started by the superuser it drops the privileges, by default it works
    as \emph{clamav}. \emph{freshclam} downloads the database from the Clam
    AntiVirus homepage and checks its consistency using MD5 sum.
    process for Clam AntiVirus.
    \textbf{Run \emph{freshclam} (as root) without any parameters to check
    is it working correctly}. If everything is OK, create the log file in
    /var/log owned by \emph{clamav}:
    \begin{verbatim}
	# touch /var/log/clam-update.log
	# chmod 600 /var/log/clam-update.log
	# chown clamav /var/log/clam-update.log
    \end{verbatim}
    Now you can run \emph{freshclam} as a daemon:
    \begin{verbatim}
	# freshclam -d -c 2 -l /var/log/clam-update.log
    \end{verbatim}
    It will check for a new database 2 times a day. Please add the above line
    to your startup scripts. The other way is to use the \emph{cron} daemon.
    You have to add a similar line to the crontab of \textbf{root} or
    \textbf{clamav}:
    {\small
    \begin{verbatim}
0 8	* * *	/usr/local/bin/freshclam --quiet -l /var/log/clam-update.log
    \end{verbatim}}
    \noindent It will check for a new database daily at 8 am. You may need
    to setup the proxy support on your system. You should set the environment
    variable \emph{\$http\_proxy}, eg.
    \begin{verbatim}
	export http_proxy="my.proxy.server:8080"
    \end{verbatim}
    There is also \emph{--http-proxy} and \emph{--proxy-user} option available.

    \subsection{FreshClam: Mirrors and mirrors.txt}
    The main server is \url{clamav.elektrapro.com} and there are the following
    mirrors available:
    \begin{itemize}
	\item \url{clamav.ozforces.com} - database mirror updated manually
	\item \url{clamav.essentkabel.com} - full mirror of the main site
	      updated automatically
	\item \url{clamav.linux-sxs.org} - database mirror (rsync from ozforces)
    \end{itemize}
    In the database directory you will find \emph{mirror.txt} file, which
    freshclam reads each time it tries to download the new database. It
    uses the first server from the file and switches to another one (and
    remembers that position for some time) when the previous is not available.
    You can modify that file if you will find some mirror faster, however
    this is not recommended.


    \section{Usage}

    \subsection{Clam daemon}\label{clamd}
    \emph{clamd} is a fully multi-threaded daemon, based on \emph{libclamav}.
    It's able to work in one of the two modes, using:
    \begin{itemize}
	\item Unix (local) sockets
	\item TCP sockets
    \end{itemize}
    The daemon is configured by the \emph{clamav.conf} file. You will find
    a description of all the options in the \textbf{clamav.conf(5)} manual.
    \emph{clamd} recognizes the following commands:
    \begin{itemize}
	\item \textbf{PING}\\
	    Check server's state. It should reply with "PONG".
	\item \textbf{VERSION}\\
	    Print the version information.
	\item \textbf{RELOAD}\\
	    Reload the databases.
	\item \textbf{QUIT}\\
	    Perform a clean exit.
	\item \textbf{SCAN file/directory}
	    Scan a file or directory (recursively) with archive support. A
	    full path is required.
	\item \textbf{RAWSCAN file/directory}
	    Scan a file or directory (recursively) with archive support
	    disabled. A full path is required.
	\item \textbf{CONTSCAN file/directory}
	    Scan a file or directory (recursively) with archive
	    support enabled and  continue  scanning  even  when
	    virus was found. A full path is required.
	\item \textbf{STREAM}
	    Scan stream - on this command clamd will return "PORT number"
	    and you can connect to that port and send a data to scan.

    \end{itemize}
    Internal threads (except clamuko) are ignoring all external signals.
    The main thread handles \emph{SIGTERM} and \emph{SIGINT} signals
    and performs a proper exit when one of them is caught.

    \subsection{Clamuko}\label{clamuko}
    Clamuko is a special thread in \emph{clamd}, that performs on-access
    scanning under Linux. It was implemented as a thread in clamd because
    of Dazuko implementation. Client (clamuko) - server (clamd) model is
    currently not supported by Dazuko. There are some benefits from
    current implementation - clamuko is sharing the database with clamd,
    and it's updated with the RELOAD command. \textbf{You must obey the
    following principles when using clamuko:}
    \begin{itemize}
	\item Always stop the daemon cleanly, with QUIT command or
	      SIGTERM signal. In other case, you can lose an access
	      to the protected files until the system is restarted.
	\item Never protect the directory your mail-scanner software
	      uses for attachments unpacking. Access to all infected
	      files will be blocked, and the scanner (even clamd)
	      won't be able to detect a virus. Infected mail will be
	      delivered.
    \end{itemize}
    You need to enable clamuko in \emph{clamav.conf}. To protect directory
    /home, please use the option:
    \begin{verbatim}
	ClamukoIncludePath /home
    \end{verbatim}
    To protect the whole system:
    \begin{verbatim}
	ClamukoIncludePath /
	ClamukoExcludePath /proc
	ClamukoExcludePath /tempdir/of/mail/scanner
    \end{verbatim}
    You can use clamuko to protect file access on Samba/Netatalk. NFS
    is not supported (Dazuko doesn't intercept NFS access calls). Another
    idea - you can build a database containing a signatures of the popular
    exploits, it will protect you against script-kiddies.

    \subsection{Archives and compressed files}
    Clam AntiVirus depends on LibClamAV. It has built-in support for the
    following formats:
    \begin{itemize}
	\item Zip
	\item Gzip
	\item RAR (2.0 only)
    \end{itemize}
    Archive files are detected by checking a magic strings.\footnote{Just like
    the file(1) command.}
    You need the zlib library for the Zip/Gzip support. Zip archives are
    accessed with the zziplib library by Guido Draheim and Tomi Ollila.
    RAR support is based on the UniquE RAR File Library by Christian Scheurer
    and Johannes Winkelmann. Both of them are included and slightly modified
    in the clamav sources. Unrarlib supports RAR 2.0 archives only and
    according to Christian the new format (introduced in WinRAR 3.0) won't
    be supported. 

    The daemon scans archives supported by libclamav only. Clamscan tries
    to scan an archive with built-in code, but when it fails it's able
    to switch to the external unpacker:
    \begin{verbatim}
	$ clamscan --unrar rarfail.rar
	/home/zolw/Clam/test/rarfail.rar: RAR module failure.

	UNRAR 3.00 freeware      Copyright (c) 1993-2002 Eugene Roshal


	Extracting from /home/zolw/Clam/test/rarfail.rar

	Extracting  test1                                           OK 
	All OK
	/tmp/44694f5b2665d2f4/test1: ClamAV-Test-Signature FOUND
	/home/zolw/Clam/test/rarfail.rar: Infected Archive FOUND
    \end{verbatim}

    clamscan supports many popular compressors - it uses external programs
    for each format. \textbf{If the scanner runs with superuser privileges
    unpackers are executed with \emph{clamav} privileges, which makes the
    process far more secure.} It also makes sure, that \emph{clamav} user
    has read access to all scanned compressed files.  \textbf{You should have
    enabled recursive scanning with the \emph{-r} option (\emph{--recursive}),
    if you want to scan the whole content of the archive (with subdirectories)},
    also all archives in archives will be recursively scanned - just everything.    If files in archives are virus free the archive itself is scanned - just
    for prevention (it may not be an archive). Please look at the options
    below, each option has an optional argument - the absolute path to unpacker.
    If it can't be found in \emph{\$PATH} please supply it. \emph{Because Clam
    AntiVirus uses the standard GNU options format, the long options with
    optional arguments, you \textbf{must} remember about the $=$ between option
    and argument.  So the proper way to supply the optional arguments is for
    example --unzip=/path/to/unzip.} \\[5pt]

    \noindent
    \up{unzip}{You probably don't need this option, because Zip is supported
	       by libclamav. But if libclamav will fail to unzip some file,
	       it may be useful.
	       clamscan was tested with \emph{UnZip 5.41 of 16 April 2000,
	       by Info-ZIP}.}
    \up{unrar}{Tested with \emph{UNRAR 3.00 freeware}.}
    \up{unace}{It uses options supported by \emph{UNACE v1.2 public version},
	       not tested, but should work.}
    \up{arj}{Tested with \emph{arj 3.10b}.}
    \up{zoo}{Tested with \emph{zoo 2.1}.}
    \up{lha}{Tested with \emph{LHa for Unix V 1.14e}.}
    \up{jar}{CA uses \emph{unzip} for .jar files. Tested with \emph{UnZip 5.41
	     of 16 April 2000, by Info-ZIP}.}
    \up{tar}{This option supports non-compressed archives. Tested with
	     \emph{GNU tar 1.13.17}.}
    \up{deb}{This option supports debian binary packages. Tested with
	     \emph{GNU ar\\ 2.12.90.0.14}. Implies --tgz , but doesn't conflict
	     with --tgz=FULLPATH.}
    \up{tgz}{This option supports .tar.gz and .tgz files. You need \emph{GNU
	     tar}, on non-Linux system you probably have it as \emph{gtar}
	     and if this is in \emph{\$PATH} just use \emph{--tgz=gtar} or
	     supply the full path to this command as an argument.}


    \subsection{Output format}
    \emph{clamd} uses clamscan compatible (see below) output format.
    \begin{verbatim}
	zolw@Wierszokleta:~$ telnet localhost 3310
	Trying 127.0.0.1...
	Connected to localhost.
	Escape character is '^]'.
	SCAN /home/zolw/infected
	/home/zolw/infected/sobre.com: W32/Magistr.B FOUND
	Connection closed by foreign host.
    \end{verbatim}
    In \textbf{SCAN} mode it closes the connection when first virus is found.
    In the case of archives, the output is exactly the same as with normal
    files:
    \begin{verbatim}
	SCAN /home/zolw/Clam/test/test2.zip
	/home/zolw/Clam/test/test2.zip: ClamAV-Test-Signature FOUND
    \end{verbatim}
    \textbf{CONTSCAN} displays all infected files found.\\
    Error messages are printed in the following format:
    \begin{verbatim}
	SCAN /no/such/file
	/no/such/file: Can't stat() the file ERROR
    \end{verbatim}
    and they can be easily parsed.

    \emph{clamscan} writes all messages (only help is written to \textbf{stdout}    by default) to \textbf{stderr}. In some situations
    you may want to redirect it to \textbf{stdout} with \emph{--stdout}.
    \emph{stdout} in contrast to \emph{stderr} is buffered, that's why
    \emph{clamscan} flushes this buffer after each message, to prevent
    the creation of trashes on the output. During scanning it writes something
    like this:
    \begin{verbatim}
	/TEST/test: OK
	/TEST/Makefile: OK
	/TEST/getopt.c: OK
	/TEST/virfile: Phantom #1 FOUND
    \end{verbatim}
    When a virus is found, its name is printed between \emph{filename:} and
    \emph{FOUND}.
    If a virus is found in an archive scanned with an external unpacker
    it's noticed with \emph{Infected Archive}. "Infected Archives" are not
    counted as infected files - just files in them are. Please note the
    difference between an internal unarchiving - because it's realized
    transparently by the libclamav, clamscan doesn't even know the file is
    an archive.

    \subsection{Signature Tool}
    \emph{sigtool} automates signature creation. If you have an infected file,
    which isn't detected by ClamAV, but it is by another anti-virus scanner
    working in the console, you can create the signature easily.
    \emph{Example of usage:}
    Create a random file and put the \textbf{test1} file content into it. We
    will use \emph{clamscan} to generate the signature, it's just an example.
    Scan it with \emph{clamscan --stdout testfile}, the output is
    \begin{verbatim}
	testfile: ClamAV-Test-Signature FOUND

	----------- SCAN SUMMARY -----------
	Known viruses: 7734
	Scanned directories: 0
	Scanned files: 1
	Data scanned: 0.95 MB
	Infected files: 1
	I/O buffer size: 131072 bytes
	Time: 0.245 sec (0 m 0 s)
    \end{verbatim}
    The unique string in this output is "ClamAV-Test-Signature". Run
    \emph{sigtool} with the following parameters:
    \begin{verbatim}
	$ sigtool -c "clamscan --stdout" -f testfile -s "ClamAV-Test"
    \end{verbatim}
    The program will concatenate arguments for \emph{-c (--command)} and
    \emph{-f (--file)}, that's why the scanner's options must be given in the
    proper order. At the end it will generate a file \emph{testfile.sig},
    which should contain 100 bytes in our example. It contains the proper
    signature.
    \begin{verbatim}
	...
	...
	Detected at 12103, moving backward.
	Detected at 11983, moving backward.
	Detected at 11923, moving backward.
	Not detected, increasing pos 11893 -> 11923
	Detected at 11923, moving backward.
	Not detected, increasing pos 11908 -> 11923
	Detected at 11923, moving backward.
	Not detected, increasing pos 11915 -> 11923
	Detected at 11923, moving backward.
	Detected at 11919, moving backward.
	Detected at 11917, moving backward.
	Detected at 11916, moving backward.
	Starting precise loop
	 *** Found signature end at 11916

	The scanner was executed 46 times.
	Signature length is 50, so length of hex string should be 100
	Saving signature in testfile.sig file.
    \end{verbatim}

    \section{Compatible software}
    The following software supports ClamAV. It's specified which elements
    are supported, please note that if the program doesn't
    support clamd you can use clamdscan instead of clamscan.

    \subsection{clamav-milter}
    \textbf{Location:} included in clamav package\\
    \textbf{Supports:} clamd\\[4pt]
    clamav-milter by Nigel Horne is a very fast email scanner designed for
    sendmail. It's entirely written in C and uses ClamAV's internal mail
    scanner (also written by Nigel).\\[4pt]
    \textbf{Installation:}\\
    You need libmilter development files. Configure ClamAV with
    \begin{verbatim}
    $ ./configure --enable-milter
    \end{verbatim}
    and recompile. The program will be installed in
    /usr/local/sbin/clamav-milter. The following instructions were adopted
    from Nigel's INSTALL file: add to /etc/mail/sendmail.mc:
    \begin{verbatim}
INPUT_MAIL_FILTER(`clmilter',`S=local:/var/run/clmilter.sock,
F=, T=S:4m;R:4m')dnl
define(`confINPUT_MAIL_FILTERS', `clmilter')
    \end{verbatim}
    Check entries in clamav.conf of the form:
    \begin{verbatim}
	LocalSocket /var/run/clamd.sock
	ScanMail
	StreamSaveToDisk
    \end{verbatim}
    Start clamav-milter:
    \begin{verbatim}
	/usr/local/sbin/clamav-milter -blo /var/run/clmilter.sock
    \end{verbatim}
    and restart sendmail.

    \subsection{mod\_clamav}
    \textbf{Location:} \url{http://software.othello.ch/mod\_clamav}\\
    \textbf{Supports:} libclamav, clamd\\[4pt]
    mod\_clamav is an Apache virus scanning filter. It was written
    and is currently maintained by Andreas M�ller. The project is very well
    documented and the installation is quite easy.

    \subsection{TrashScan}
    \textbf{Location:} clamav-sources/support/trashscan\\
    \textbf{Supports:} clamscan\\[4pt]

    This is a procmail based scanner from Trashware and it's extremely
    easy to setup, however this is for single users only and not as
    efficient as MTA based scanners.

    \subsection{AMaViS - "Next Generation"}
    \textbf{Location:} \url{http://sourceforge.net/projects/amavis}\\
    \textbf{Supports:} clamscan\\[4pt]
    AMaViS-ng is a rewritten, more modular version of amavis-perl/amavisd,
    developed by Hilko Bengen. Home site:\\[4pt]
    \textbf{Installation:}\\

    \noindent Please download the newest version (at least 0.1.4).
    After installation (which is quite easy), please uncomment the following
    line in amavis.conf:
    \begin{verbatim}
	virus-scanner = CLAM
    \end{verbatim}
    and eventually change the path to clamscan in the \emph{[CLAM]} section:
    \begin{verbatim}
	[CLAM]

	clamscan = /usr/local/bin/clamscan
    \end{verbatim}

    \subsection{amavisd-new}
    \textbf{Location:} \url{http://www.ijs.si/software/amavisd}\\
    \textbf{Supports:} clamd, clamscan\\[4pt]
    amavisd-new is a rewritten version of amavis maintained by
    Mark Martinec.\\[4pt]
    \textbf{Installation:}\\
    clamscan is enabled automatically if clamscan binary is found
    at amavisd-new starup time. clamd is activated by uncommenting
    its entry in the @av\_scanners list, file /etc/amavisd.conf.

    \subsection{Qmail-Scanner}
    \textbf{Location:} \url{http://qmail-scanner.sf.net}\\
    \textbf{Supports:} clamscan\\[4pt]
    You must increase softlimit value or wait for a daemon support.

    \subsection{Sagator}
    \textbf{Location:} \url{http://www.salstar.sk/sagator}\\
    \textbf{Supports:} clamscan, clamd, libclamav\\[4pt]
    This program is an email antivirus/antispam gateway. It is an interface
    to the postfix (or any other smtpd), which runs antivirus
    and/or spamchecker. Its modular architecture can use any
    combination of antivirus/spamchecker according to configuration.

    \subsection{ClamdMail}
    \textbf{Location:} \url{http://clamdmail.sf.net}\\
    \textbf{Supports:} clamd\\[4pt]
    A mail processing client for ClamAV. Small, fast and easy to install.

    \subsection{BlackHole}
    \textbf{Location:} \url{http://www.groovy.org/blackhole.shtml}\\
    \textbf{Supports:} clamscan, clamd\\[4pt]
    BlackHole is an advanced spam / virus filter for Qmail, Postfix, Sendmail,
    Exim and Courier written by Chris Kennedy. This tool is for advanced
    administrators (installation is hard).

    \subsection{MailScanner}
    \textbf{Location:} \url{http://www.mailscanner.info}\\
    \textbf{Supports:} clamscan\\[4pt]
    MailScanner scans all e-mail for viruses, spam and attacks against
    security vulnerabilities. It is not tied to any particular virus
    scanner, but can be used with any combination of 14 different virus
    scanners, allowing sites to choose the "best of breed" virus scanner.

    \subsection{MIMEDefang}
    \textbf{Location:} \url{http://www.roaringpenguin.com/mimedefang}\\
    \textbf{Supports:} clamscan, clamd\\[4pt]
    This is an efficient mail scanner for Sendmail/milter.

    \subsection{Exiscan}
    \textbf{Location:} \url{http://duncanthrax.net/exiscan}\\
    \textbf{Supports:} clamscan, clamd\\[4pt]
    exiscan is a patch against exim version 4, providing support for content
    scanning in email messages received by exim. Four different scanning
    facilities are supported: antivirus, antispam, regular expressions, and
    file extensions.

    \section{LibClamAV}
    libclamav may be used to add a virus protection into your software.
    The library is thread-safe, automatically recognizes and scans an
    archives. Scanning is very fast - in most cases it won't be noticeable.

    \subsection{API}
    Each program using libclamav must include \emph{clamav.h} header file:
    \begin{verbatim}
	#include <clamav.h>
    \end{verbatim}
    The first step is an engine initialization. There are three functions
    available:
    \begin{verbatim}
	int cl_loaddb(const char *filename, struct cl_node **root,
	int *virnum);

	int cl_loaddbdir(const char *dirname, struct cl_node **root,
	int *virnum);

	char *cl_retdbdir(void);

    \end{verbatim}
    \emph{cl\_loaddb()} loads one database per time, \emph{cl\_loaddbdir()}
    loads all \emph{.db} and {.db2} files from the directory \emph{dirname}.
    \emph{cl\_retdbdir()} returns hardcoded database directory path.
    The database will be saved under \emph{root} and the number of the loaded
    signatures will be \textbf{added} to \emph{virnum}. Pointer to the tree
    structure (trie, see \ref{engine}) must initially point to the NULL. If you
    don't want to save the number of signatures loaded pass the NULL as the
    third argument. \emph{cl\_loaddb} functions return 0 on success and
    other value on failure.
    \begin{verbatim}
	    struct cl_node *root = NULL;
	    int ret;

	ret = cl_loaddbdir(cl_retdbdir(), &root, NULL);
    \end{verbatim}
    There's elegant way to print libclamav's error codes:
    \begin{verbatim}
	char *cl_perror(int clerror);
    \end{verbatim}
    \emph{cl\_perror()} returns a (statically allocated) string describing
    \emph{clerror} code:
    \begin{verbatim}
	if(ret) {
	    printf("cl_loaddbdir() error: %s\n", cl_perror(ret));
	    exit(1);
	}
    \end{verbatim}
    When database is loaded, you must create the proper trie with:
    \begin{verbatim}
	void cl_buildtrie(struct cl_node *root);
    \end{verbatim}
    In our example:
    \begin{verbatim}
	cl_buildtrie(root);
    \end{verbatim}
    OK, now you can scan a buffer, descriptor or file with:
    \begin{verbatim}
	int cl_scanbuff(const char *buffer, unsigned int length,
	char **virname, const struct cl_node *root);

	int cl_scandesc(int desc, char **virname, unsigned long int
	*scanned, const struct cl_node *root, const struct cl_limits
	*limits, int options);

	int cl_scanfile(const char *filename, char **virname,
	unsigned long int *scanned, const struct cl_node *root,
	const struct cl_limits *limits, int options);
    \end{verbatim}
    All the functions save a virus name address under \emph{virname} pointer.
    \emph{virname} points to the name in the trie structure, thus it can't be
    released directly. \emph{cl\_scandesc()} and \emph{cl\_scanfile()} can
    increase \emph{scanned} value in CL\_COUNT\_PRECISION units. They also
    support archive limits:
    \begin{verbatim}
	struct cl_limits {
	    int maxreclevel;
	    int maxfiles;
	    long int maxfilesize;
	};
    \end{verbatim}
    The last argument configures scan engine. Currently it supports
    \textbf{CL\_ARCHIVE} (enables archive scanning), \textbf{CL\_RAW}
    (disables archive scanning) and \textbf{CL\_MAIL} (enables mbox
    and Maildir scanning).
    The functions return 0 (\textbf{CL\_CLEAN}) when no virus is found,
    \textbf{CL\_VIRUS} when virus is found and other value on failure.
    \begin{verbatim}
	    struct cl_limits limits;
	    char *virname;

	/* maximal number of files in archive */;
	limits.maxfiles = 100
	/* maximal archived file size == 10 MB */
	limits.maxfilesize = 10 * 1048576;
	/* maximal recursion level */
	limits.maxreclevel = 8;


	if((ret = cl_scanfile("/home/zolw/test", &virname, NULL, root,
	&limits, CL_ARCHIVE)) == CL_VIRUS) {
	    printf("Detected %s virus.\n", virname);
	} else {
	    printf("No virus detected.\n");
	    if(ret != CL_CLEAN)
	        printf("Error: %s\n", cl_perror(ret));
	}
    \end{verbatim}
    When you don't need to scan more files, the trie should be released
    with:
    \begin{verbatim}
	void cl_freetrie(struct cl_node *root);
    \end{verbatim}
    You will find some examples in clamav sources. Each program using
    libclamav must be linked against it:
    \begin{verbatim}
	gcc -Wall ex1.c -o ex1 -lclamav
    \end{verbatim}
    Enjoy !


    \section{Problem solving}

    \subsection{Return codes}
    Return codes are very useful, especially in system scripts. You may
    check the return code from \emph{clamscan}, by running the following
    command directly after the scanner exits:
    \begin{verbatim}
	$ echo $?
    \end{verbatim}
    Here is a list of return codes from \emph{clamscan}:\\[4pt]

    \noindent
    \rc{0}{No virus was found.}
    \rc{1}{Virus(es) detected.}
    \rc{40}{Unknown option was passed to \emph{clamscan}. Please check
	    \emph{clamscan --help} or manual page for available options.}
    \rc{50}{Problem with initialization of virus database. Probably
	    it doesn't exist in the default place or wrong file was passed
	    to \emph{--database}.}
    \rc{51}{Wrong number of threads was passed to \emph{--threads}. It
	    must be a natural number $ \ge 0$.} 
    \rc{52}{Not supported file type. Scanner supports regular files,
	    directories and symlinks.}
    \rc{53}{Can't open directory.}
    % FIXME: ^5
    \rc{54}{Can't open file.$^5$}
    \rc{55}{Error reading file. Probably the medium you are reading is broken.
	    \footnote{Only in one-file mode (in recursive mode those errors are
	    ignored)}}
    \rc{56}{Can't stat input file or directory. File / directory you want to
	    scan doesn't exist.}
    \rc{57}{Can't get absolute pathname of current working directory. Your
	    current pathname is longer then 200 characters. When clamscan
	    is started without a input file / directory it scans the current
	    directory. For some reasons it needs absolute pathnames, the buffer
	    is hardcoded to 200 characters and that should be sufficient.}
    \rc{58}{I/O error. Please check the filesystem.}
    \rc{59}{Can't get information about current user (running clamscan).}
    \rc{60}{Can't get information about user \emph{clamav}. User \emph{clamav}
	    (default unprivileged user) doesn't exist in /etc/passwd.}
    \rc{61}{Can't fork. Can't create new process, please check your limits.}
    \rc{63}{Can't create temporary file or directory. Please check permissions.}
    \rc{64}{Can't write to temporary directory. Please specify another one.}
    \rc{70}{Can't allocate and clear memory. This is a critical error, please
	    check your system.}
    \rc{71}{Can't allocate memory. Look above.}


    \section{Technicals}

    \subsection{Security}
    Clam AntiVirus cares about security. Dangerous operations in clamscan
    (such as extracting, temporary file creation, unlink() operations) are
    executed with \emph{clamav} privileges. \textbf{But there are no programs
    without bugs.} This is a young project and everything is possible. In some
    places it uses the \emph{snprintf()} function, some older systems (C
    libraries) however the buffer length in this function isn't checked. This
    example shows, that you should check your system first. Never set SUID/SGID
    bits on Clam AntiVirus executables. If the SUID bit is set and
    \emph{clamscan} is owned by root, every file on the system may be modified
    with the \emph{--log} option. Normal users may use \emph{clamscan} to scan
    their files, other files shouldn't interest them. Clam AntiVirus Daemon
    was written with security in mind - it doesn't allow external unpackers
    (uses only libclamav unarchivers) and contains some additional protections.


    \subsection{Scan engine}\label{engine}
    New versions of Clam AntiVirus are using a mutation of Aho-Corasick
    pattern matching algorithm. This algorithm uses a finite state pattern
    matching automaton \cite{clr}. The algorithm itself is a generalization of
    the Knuth-Morris-Pratt algorithm. Please look at \emph{matcher.h} for data
    type definitions. The automaton is represented by the trie. Trie is
    a rooted tree with some specific properties \cite{acwww}. Each node
    of the trie represents some state of the automaton. In the implementation,
    the node is defined as following:
    \begin{verbatim}
	struct node {
	    int islast;
	    struct patt *list;
	    int maxpatlen;
	    struct node *next[NUM_CHILDS], *trans[NUM_CHILDS], *fail;
	};
    \end{verbatim}
    [To be continued...]

    \section{Credits}
    In alphabetical order:
    \begin{itemize}
	\item AIX PDSLIB, University of California at Los Angeles \\
	      \url{http://aixpdslib.seas.ucla.edu} - binary packages for AIX
	\item Kamil Andrusz \email{<wizz(.at.)mniam.net>} - OpenBSD support
	      patch
	\item Jean-Edouard BABIN \email{<Jeb(.at.)jeb.com.fr>} - NetBSD support;
	      made his NetBSD box available to me.
	\item Marc Baudoin \email{<babafou(.at.)babafou.eu.org>} - NetBSD
	      testing
	\item Hilko Bengen \email{<bengen(.at.)vdst-ka.inka.de>} - support for
	      Clam AntiVirus in his AMaViS - "Next Generation"
	\item Patrick Bihan-Faou \email{<patrick(.at.)mindstep.com>} - support
	      for --with-user/group in the configure script.
	\item Eric I. Lopez Carreon \email{<elopezc(.at.)technitrade.com>} -
	      Spanish\\ "Sendmail + AMaViS + ClamAV Installation" how-to
	\item Nicholas Chua \email{<nicholas(.at.)ncmbox.net>} -
	      virus signatures; clamav.org mirror administrator.
	\item Damien Curtain \email{<damien(.at.)pagefault.org>} - fix for the
	      --remove option in clamscan (it didn't work with internal
	      archivers); implementation of the --move option in clamscan,
	      mirroring support in freshclam.
	\item Krisztian Czako \email{<slapic(.at.)linux.co.hu>} - virus
	      signatures.
	\item Diego D'Ambra \email{<da@softcom.dk>} - \textbf{Database
	      developer.}
	\item Alejandro Dubrovsky \email{<s328940(.at.)student.uq.edu.au>} -
	      patch for including and excluding multiple patterns.
	\item Magnus Ekdahl \email{<magnus(.at.)debian.org>} - Debian\\
	      (\url{http://www.debian.org}) package maintainer; fixes and
	      improvements.
	\item Jason Englander \email{<jason(.at.)englanders.cc>} - bug report:
	      clamd recursive scanning of the directories on non standard
	      file systems; configure script support for id checking.
	      \textbf{Database developer.}
	\item Oden Eriksson \email{<oden.eriksson(.at.)kvikkjokk.net>} -
	      Mandrake\\ package maintainer.
	\item Edison Figueira Junior \email{<edison(.at.)brc.com.br>} - money
	      donation.
	\item David Ford \email{<david+cert(.at.)blue-labs.org>} - gcc 3.x
	      support fix.
	\item Piotr Gackiewicz \email{<gacek(.at.)intertele.pl>} - bug report:
	      clamd\\ THREXIT bug
	\item Nick Gazaloff \email{<nick(.at.)sbin.org>} - socket descriptors
	      leak fix in clamd.
	\item Wieslaw Glod \email{<wkg(.at.)x2.pl>} - bug report: FreeBSD
	      compile problem in 0.22.
	\item Matthew A. Grant \email{<grantma(.at.)anathoth.gen.nz>} -
	      OpenAntiVirus Update script (\emph{oav-update})
	\item Hrvoje Habjanic \email{<hrvoje.habjanic(.at.)zg.hinet.hr>} -
	      syslog support patch for clamd; virus provider.
	\item Michal Hajduczenia \email{<michalis(.at.)mat.uni.torun.pl>} - Clam
	      title logo.
	\item Paul Hoadley \email{<paulh(.at.)logixsquad.net} - "Installing
	      qmail-scanner, Clam AntiVirus and SpamAssassin under FreeBSD"
	      how-to.
	\item Thomas W. Holt Jr. \email{<twh(.at.)cohesive.net>} -
	      information about ClamAV compiling on Solaris 2.6 and Cobalt
	      MIPS boxes.
	\item Douglas J Hunley \email{<doug(.at.)hunley.homeip.net>} -\\
	      clamav.linux-sxs.org mirror, ideas.
	\item Kurt Huwig \email{<kurt(.at.)iku-netz.de>} - smart suggestions,
	      ScannerDaemon (OpenAntiVirus) author.
	\item Dave Jones \email{<dave(.at.)kalkbay.co.za>} - bug report: problem
	      in option parser.
	\item Kazuhiko \email{<kazuhiko(.at.)fdiary.net>} - Qmail-Scanner 0.12
	      support patch.
	\item Robbert Kouprie \email{<robbert(.at.)exx.nl>} - patch for unrarlib
	      buffer overflow.
	\item Henk Kuipers \email{<henk(.at.)opensourcesolutions.nl>} - bug
	      report:\\ 0.50 compile problem.
	\item Nigel Kukard \email{<nkukard(.at.)lbsd.net>} - virus signatures.
	\item Dr Andrzej Kurpiel \email{<akurpiel(.at.)mat.uni.torun.pl>} -
	      choice of\\ this project from my list.
	\item Dennis Leeuw \email{<dleeuw(.at.)made-it.com>} - \emph{"Debian
	      GNU/Linux Mail Server"} how-to, \textbf{corrections of this
	      document}.
	\item Free Oscar \email{<freeoscar(.at.)wp.pl>} - hex2str() enhancement
	\item Martin Lesser \email{<admin-debian(.at.)bettercom.de>} - patch for
	      the\\ http-proxy problem in 0.51.
	\item Peter N Lewis \email{<peter(.at.)stairways.com.au>} - Mac OS X
	      data type problem bugfix.
	\item Mike Loewen \email{<mloewen(.at.)sturgeon.cac.psu.edu>} - bug
	      report:\\ clamscan 0.24 compile error on Solaris 8; various
	      Solaris and AIX tips.
	\item Stefan Martig \email{<sm(.at.)officeco.ch>} - bug report:
	      /proc/cpuinfo problem analysis on Linux/Alpha, providing me with
	      access to the Linux/Alpha system.
	\item Brian May \email{<bam(.at.)debian.org>} - bug report: clamd
	      writing to an undefined file.
	\item Ken McKittrick \email{<klmac(.at.)usadatanet.com>} - intensive
	      FreeBSD\\ testing, hdd donation.
	\item Chris van Meerendonk \email{<cvm(.at.)castel.nl>} - virus
	      samples, clamav.essentkabel.com mirror.
	\item Arkadiusz Miskiewicz \email{<misiek(.at.)pld.org.pl>} - Polish(ed)
	      Linux\\ Distribution (\url{http://www.pld.org.pl}) rpm package
	      maintainer; fixes and ideas.
	\item Doug Monroe \email{<doug(.at.)planetconnect.com>} - Qmail-Scanner
	      problem analysis.
	\item Lukasz Jaroslaw Mozer \email{<lukasz(.at.)netx.waw.pl>} -
	      \url{clamav.rulez.pl} mirror.
	\item Flinn Mueller\email{flinn(.at.)activeintra.net} - OpenBSD port
	      maintainer.
	\item Hendrik Muhs \email{<Hendrik.Muhs(.at.)student.uni-magdeburg.de>}
	      -\\ pattern matcher optimization.
	\item Luca 'NERvOus' Gibelli \email{<nervous(.at.)nervous.it>} -
	      ElektraPro.com administrator.
	\item Wojciech Noworyta \email{<wnow(.at.)konarski.edu.pl>} - bug
	      report: buffer overflow in clamscan's help under Windows.
	\item Joe Oaks \email{<joe.oaks(.at.)hp.com>} - HPUX support.
	\item Washington Odhiambo \email{<wash(.at.)wananchi.com>} - extensive
	      mbox\\ code testing, bug reports.
	\item Masaki Ogawa \email{<proc(.at.)mac.com>} - Mac OS X support,
	      Japanese documentation.
	\item Martijn van Oosterhout \email{<kleptog(.at.)svana.org>} - code
	      analysis and suggestions.
	\item OpenAntiVirus.org Team - initial virus database.
	\item Tomasz Papszun \email{tomek(.at.)lodz.tpsa.pl} - various bug
	      reports.
	\item Eric Parsonage \email{eric(.at.)eparsonage.com} - "Installing
	      qmail-scanner, Clam Antivirus and SpamAssassin under FreeBSD"
	      how-to.
	\item Oliver Paukstadt \email{<pstadt(.at.)stud.fh-heilbronn.de>} - bug
	      report: crash with strange Zip archives.
	\item Kristof Petr \email{<Kristof.P(.at.)fce.vutbr.cz>} - bug report:
	      socket descriptors leak in clamd; file decriptors leak in
	      clamd, clamscan and libclamav.
	\item Ed Phillips \email{<ed(.at.)UDel.Edu>} - patch for the internal
	      logger in clamd.
	\item Andreas Piesk \email{<Andreas.Piesk(.at.)heise.de>} - clamd:\\
	      ScannerDaemonOutputFormat option.
	\item Ant La Porte \email{<ant(.at.)dvere.net>} - proxy support
	      enhancement.
	\item Sergei Pronin \email{<sp(.at.)finndesign.fi>} - bug report:
	      access problems in superuser mode.
	\item Thomas Quinot \email{<thomas(.at.)cuivre.fr.eu.org>} - patch for\\
	      non-default prefix and incoherent database location
	      specification in defaults.h of clamscan and freshclam.
	\item David Sanchez \email{<dsanchez(.at.)veloxia.com>} - bug report:
	      thread deadlocking\\ in a critical error situation.
	\item Martin Schitter - bug report: libclamav crash on certain
	      zip files.
	\item Enrico Scholz
	      \email{<enrico.scholz(.at.)informatik.tu-chemnitz.de>} -
	      daemonize() enhancements.
	\item Dr Zbigniew Szewczak \email{<zssz(.at.)mat.uni.torun.pl>} - ideas,
	      suggestions and time spent on discussing some aspects of ClamAV.
	\item Gernot Tenchio \email{<g.tenchio(.at.)telco-tech.de>} - proxy
	      authorization support in freshclam.
	\item Masahiro Teramoto \email{<markun@onohara.to>} - official FreeBSD
	      port maintainer.
	\item Trashware \email{trashware(.at.)gmx.net} - TrashScan
	\item Troy Wollenslegel \email{<troy(.at.)intranet.org>} - bug report:
	      handling inaccessible directories in archives.
	\item Andoni Zubimendi \email{<andoni(.at.)lpsat.net>} - fix for
	      segmentation\\ fault in 0.12 (NULL pointer dereference).
    \end{itemize}

    \section{Authors}
    Nigel Horne \email{<njh(.at.)bandsman.co.uk>} is an active ClamAV developer
    responsible for the mbox code in libclamav and clamav-milter. I take care
    most of these things are working ;) If you have some questions, feel free
    to mail us.\\[4pt]

    \hfill Tomasz Kojm \email{<zolw(.at.)konarski.edu.pl>}

    \begin{thebibliography}{99}
	\bibitem{clr}
	    Cormen, Leiserson, Rivest: \emph{Introduction to Algorithms},
	    Chapter 34, MIT Press.
	\bibitem{acwww}
	    {\small
	    \url{http://www-sr.informatik.uni-tuebingen.de/~buehler/AC/AC.html}}:
	    Aho-Corasick algorithm description
    \end{thebibliography}
\end{document}
